\subsection{\IT{vprintf()}関数のケース}
\myindex{\CStandardLibrary!vprintf}
\myindex{\CStandardLibrary!va\_list}

多くのプログラマは、printf のようなフォーマット文字列 + 可変数の引数を取る独自のログ関数を定義します。

もう一つの一般的な例は、何らかのメッセージを出力して終了する die() 関数です。

未知数の入力引数をパックして\printf 関数に渡す何らかの方法が必要です。
しかし、どのように？

そのため、名前に\q{v}が付く関数があります。

その一つが\IT{vprintf()}です：これはフォーマット文字列と\TT{va\_list}型の変数へのポインタを取ります：

\lstinputlisting[style=customc]{\CURPATH/die.c}

詳しく調べると、\TT{va\_list}が配列へのポインタであることがわかります。
コンパイルしてみましょう：

\lstinputlisting[caption=\Optimizing MSVC 2010,style=customasmx86]{\CURPATH/die_MSVC2010_Ox_EN.asm}

私たちの関数が行うことは、単に引数へのポインタを取得し、それを\IT{vprintf()}に渡すことであり、その関数はそれを無限の引数配列として扱うことがわかります！

\lstinputlisting[caption=\Optimizing MSVC 2012 x64,style=customasmx86]{\CURPATH/die_MSVC2012_x64_Ox_EN.asm}